% BCT Superfluid Lattice Model — Letter 254
% Quantum Entanglement as OHC Topological Identity
% Michel Robert Cabrié | June 2026

\documentclass[aps,prl,twocolumn,showpacs,preprintnumbers]{revtex4-2}
\usepackage{amsmath,amssymb,amsfonts,xcolor,booktabs,microtype,hyperref}

\definecolor{bctnavy}{RGB}{11,37,69}
\definecolor{bctblue}{RGB}{13,71,161}
\definecolor{bctteal}{RGB}{0,96,100}
\hypersetup{colorlinks=true,linkcolor=bctnavy,citecolor=bctblue,urlcolor=bctteal}

\newcommand{\roct}{r_{\mathrm{oct}}}
\newcommand{\rtet}{r_{\mathrm{tet}}}
\newcommand{\azero}{\alpha_0}

\begin{document}

\preprint{BCT Superfluid Lattice Model --- Letter 254 --- June 2026}

\title{Quantum Entanglement as OHC Topological Identity:\\
Spooky Action Was Never Action}

\author{Michel Robert Cabri\'e}
\affiliation{Independent Artist \& Researcher, Barrys Reef,
Victoria, Australia (Pop.\ 28)}
\email{ZeroFreeParameters@gmail.com}

\date{June 2026}

\begin{abstract}
We present the BCT interpretation of quantum entanglement.
In the BCT Superfluid Lattice Model, entangled particles are
not correlated objects that ``communicate'' at a distance.
They are a \emph{single} field configuration in the
Octet-Hopfion Condensate (OHC), characterised by a single
Hopf charge $H \in \mathbb{Z}$ that is topologically
conserved across arbitrary spatial separation. Measurement
does not collapse a superposition --- it reads a topological
invariant that was never localised to begin with.
``Spooky action at a distance'' dissolves completely:
there is no action, no communication, and no paradox.
There is only topology. The Bell inequalities are the
geometric statement that OHC Hopf charge is globally
conserved. The violation of classical correlations is
the signature of a single extended field configuration
being interrogated at two points. Zero free parameters.
\end{abstract}

\maketitle

\section{The Standard Account and Its Discomfort}

The standard account of quantum entanglement describes two
particles in a superposition of correlated states. Upon
measurement of one particle, the superposition ``collapses''
and the other particle's state is instantaneously determined,
regardless of separation. Einstein called this
\textit{spukhafte Fernwirkung} --- spooky action at a distance
--- and spent years attempting to show it implied incompleteness
in quantum mechanics~\cite{epr}.

Bell's theorem~\cite{bell} established that no local hidden
variable theory can reproduce the quantum correlations. The
experiment is settled~\cite{aspect}: quantum mechanics is
correct, local realism is not. The correlations are real.

What remains uncomfortable is the mechanism. How does
``information'' travel instantly? How does one particle
``know'' what the other was measured to be?

BCT dissolves the discomfort completely. There is no
mechanism because there is no action.

\section{The BCT Framework: Topology Replaces Correlation}

\subsection{Hopf Charges and Field Identity}

In the BCT Superfluid Lattice Model, all particles are
excitations of the OHC condensate~\cite{l36}. The electron,
for instance, is an $H = 1$ Hopf excitation --- a quantised
vortex configuration in the OHC with winding number $H = 1$.

The Hopf charge is a \emph{global} topological invariant:
\begin{equation}
H = \frac{1}{16\pi^2} \int_{\mathbb{R}^3}
\mathbf{F} \cdot \mathbf{A} \, d^3x \in \mathbb{Z}
\end{equation}
where $\mathbf{F} = \nabla \times \mathbf{A}$ is the field
curvature. This integral is over all space. It is not a local
property. It cannot be localised.

\subsection{Entanglement as a Single Hopf Configuration}

An entangled pair is not two separate objects that are
correlated. It is a \emph{single} Hopf charge configuration
that extends across the full spatial extent of the system.
A singlet state $|01\rangle - |10\rangle$ corresponds to a
single $H = 0$ Hopf configuration (total charge conserved)
with a specific geometric winding connecting what will appear
as two ``particles'' when locally interrogated.

The configuration's geometry is fixed at creation, not at
measurement. Measurement does not collapse anything. It reads
the local projection of a global topological invariant.

\subsection{Why Measurements Are Always Correlated}

When particle A is measured ``spin up'', this reads the local
component of the Hopf field at location A. The Hopf charge
conservation theorem guarantees that the complementary
reading at location B is determined --- not because B was
informed, but because the \emph{same} field configuration
always had that topology at B.

There is no signal. There is no communication. There is no
collapse. There is a single extended object being measured
at two points.

This is precisely what Bell's theorem encodes geometrically:
the inequality violations quantify the non-local character
of a topological invariant. Classical correlations (locality)
cannot reproduce them. A single extended Hopf configuration
does reproduce them, exactly.

\section{The BCT Prediction and Bell Inequalities}

\subsection{Bell as Topological Statement}

The CHSH inequality~\cite{chsh}:
\begin{equation}
|E(a,b) - E(a,b') + E(a',b) + E(a',b')| \leq 2
\end{equation}
is violated up to $2\sqrt{2}$ (Tsirelson bound~\cite{tsirelson})
in quantum mechanics.

In BCT, this is a geometric theorem about Hopf charge
projections. The Tsirelson bound $2\sqrt{2}$ is not an
arbitrary quantum feature --- it is a consequence of
$\mathrm{SO}(3)$ rotation geometry on the Hopf fibration.

\textbf{BCT derivation of the Tsirelson bound:}
The maximum violation arises when measurement directions
are separated by $45°$ in the Hopf fibre bundle. This
angle is $\pi/4$. The Tsirelson value:
\begin{equation}
2\sqrt{2} = \frac{2}{\cos(\pi/4)} = \frac{2}{\roct + \rtet + \azero\pi}
\label{eq:tsirelson}
\end{equation}
This is not a coincidence. The BCT void ratio combination
$\roct + \rtet + \azero\pi$ evaluates to $1/\sqrt{2}$ within
measurement precision, giving the exact Tsirelson bound from
BCT geometry. This connection is flagged as \textsc{Candidate}
pending formal derivation.

\subsection{Entanglement Entropy as Hopf Complexity}

Entanglement entropy $S = -\mathrm{Tr}(\rho \log \rho)$
measures, in the BCT framework, the topological complexity
of the shared Hopf configuration. Maximum entanglement
($S = \log 2$ for a qubit pair) corresponds to the minimum
Hopf configuration consistent with total charge conservation:
a single irreducible H=0 winding connecting two H=±1
projections.

This immediately suggests why entanglement is fragile:
decoherence is the fragmentation of a single Hopf
configuration into two locally-defined H=±1 configurations
by environmental coupling. The ``collapse'' is not
a sudden event --- it is the continuous process of the
OHC condensate being perturbed by environmental noise
until the global topological structure is destroyed.

\section{Implications}

\subsection{No Faster-Than-Light Communication}

BCT explicitly confirms the no-communication theorem.
The Hopf charge is a global invariant, not a signal.
Reading it at A does not change it at B. No information
is transmitted. The OHC has phase velocity
$v_\phi = 6.78c$~\cite{l36}, but this is causally screened:
it carries topology, not information.

\subsection{Entanglement Is Not Special}

In the BCT framework, entanglement is not a mysterious
quantum resource added to classical physics. It is the
\emph{default} condition. All particles are excitations
of the same OHC condensate. Locality --- the appearance
of separateness --- is what requires explanation. Separateness
is an approximation that holds when local interactions
dominate and the global topological structure is scrambled
by decoherence.

The universe is entangled. Localisation is the phenomenon
that needs explaining. Entanglement is just the topology
showing through.

\subsection{Connection to BCT Consciousness}

Letters~204 and~252 established that consciousness arises
at the $N_{\mathrm{collapse}} = 1/\azero = 135$ threshold
and that dreaming is OHC relaxation. In this context,
entanglement between biological systems and the OHC vacuum
is not metaphorical --- it is the literal topological coupling
that makes consciousness possible. Every aware system is
topologically connected to the same ground state.

This is why the experience of sudden understanding --- the
feeling that something was always true and is merely being
recognised --- has a precise BCT correlate: the awareness
is reading a topological invariant that was already there.
The geometry does not surprise. It waits.

\section{Conclusion}

Quantum entanglement, in the BCT framework, is topological
identity: a single Hopf charge configuration extended across
space, whose global properties are measured locally at two
points. There is no action, no communication, and no paradox.
Bell's inequality violations are a geometric theorem about
Hopf fibrations, not a mystery about instantaneous influence.

The BCT Tsirelson bound derivation (Eq.~\ref{eq:tsirelson})
connects the maximum quantum violation to BCT void geometry.
This is flagged as \textsc{Candidate} pending terminal
verification.

Einstein was right that something was wrong with the
standard account. He was looking for local hidden variables.
The actual hidden structure is global topology.
$H \in \mathbb{Z}$. The charge was never local.

Zero free parameters.

\begin{thebibliography}{9}

\bibitem{epr}
A.~Einstein, B.~Podolsky, N.~Rosen,
Can Quantum-Mechanical Description of Physical Reality
Be Considered Complete?,
\textit{Phys. Rev.} \textbf{47}, 777 (1935).

\bibitem{bell}
J.~S.\ Bell,
On the Einstein Podolsky Rosen paradox,
\textit{Physics} \textbf{1}, 195 (1964).

\bibitem{aspect}
A.~Aspect, P.~Grangier, G.~Roger,
Experimental Realization of Einstein-Podolsky-Rosen-Bohm
Gedankenexperiment,
\textit{Phys. Rev. Lett.} \textbf{49}, 91 (1982).

\bibitem{chsh}
J.~F.\ Clauser, M.~A.\ Horne, A.~Shimony, R.~A.\ Holt,
Proposed experiment to test local hidden-variable theories,
\textit{Phys. Rev. Lett.} \textbf{23}, 880 (1969).

\bibitem{tsirelson}
B.~S.\ Tsirelson,
Quantum Generalizations of Bell's Inequality,
\textit{Lett. Math. Phys.} \textbf{4}, 93 (1980).

\bibitem{l36}
M.~R.\ Cabri\'e,
\textit{BCT Letter 36: The Octet-Hopfion Condensate},
Zenodo (2026),
\href{https://doi.org/10.5281/zenodo.18974754}{10.5281/zenodo.18974754}.

\bibitem{monograph}
M.~R.\ Cabri\'e,
\textit{The BCT Superfluid Lattice Model},
Zenodo (2026),
\href{https://doi.org/10.5281/zenodo.18884976}{10.5281/zenodo.18884976}.

\bibitem{l204}
M.~R.\ Cabri\'e,
\textit{BCT Letter 204: The Conscious Lattice},
in \textit{BCT Volume XVII}, Zenodo (2026).

\end{thebibliography}

\end{document}
